The massless example is special. Its extra fixing can change even when the
metric and its causal curves are held fixed. Let $\mathcal L_m$ consist of
all smooth real solutions on the closed dust strip of
\[
 (\Box_g-R_g/6-m^2)\phi=0,
\]
with ordinary constant physical mass $m\ge0$, and write $F_m(b)$ for its
fixing set. Set $\Lambda^2=-\Delta_{S^3}+1$ and $T=2\pi$. With
$\psi=a\phi$ and $\lambda=m^2$, the rescaled equation is
\begin{equation}\label{eq:mass}
 \psi_{\lambda,\eta\eta}+(\Lambda^2+\lambda q(\eta))\psi_\lambda=0,
 \qquad q(\eta)=a(\eta)^2.
\end{equation}
In particular a constant addition to the rescaled potential would be a
different perturbation, corresponding to physical potential $\delta/a^2$.

\begin{proposition}[Loss of exact fixing under ordinary mass]\label{prop:mass}
Fix the dust amplitude $a_m>0$. For each $b=(\tau,Y)$ with $0<\tau<\pi$,
there are a fixed event
$x_*=(T-s,Y)$, $0<s<\tau$, and $m_*>0$ such that
\[
 x_*\in F_0(b)\cap\Isp(b),\qquad
 x_*\notin F_m(b)\quad\text{for every }0<m<m_*.
\]
The background metric is the same for all these solution classes.
\end{proposition}

\begin{proof}
Choose a nonzero smooth $v\ge0$ supported strictly outside
$\bar B(Y,\tau)$, and solve \eqref{eq:mass} with data $(0,v)$ at the
auxiliary cylinder slice $\eta=0$. The coefficient $q=a^2$ is smooth through
$0$ and $T$, so ordinary smooth Cauchy evolution and parameter dependence
apply on the extended cylinder. This is a device for constructing solutions
on the open physical strip, not a physical continuation through its
singularities. For $z=(\eta,Z)\in\Is(b)$, the initial ball
$\bar B(Z,\eta)$ lies in $B(Y,\tau)$ by the triangle inequality. Finite
propagation therefore gives $\psi_\lambda=0$ throughout $\Is(b)$ for every
$\lambda$; the zero-order term does not change this conclusion.

Put $S_0(u)=\sin(u\Lambda)/\Lambda$. The first variation
$w=\partial_\lambda\psi_\lambda|_{\lambda=0}$ has zero initial data and source
$-q(u)S_0(u)v$. Since $S_0(T-u)=-S_0(u)$, Duhamel's formula gives
\[
 w(T)=\int_0^T q(u)S_0(u)^2v\,du.
\]
The frequencies are $k=n+1$. Expanding
$q=(a_m^2/4)(3/2-2\cos u+\tfrac12\cos2u)$ gives
\[
 \int_0^T q(u)\frac{\sin^2(ku)}{k^2}\,du
 =\begin{cases}5\pi a_m^2/16,&k=1,\\
                 3\pi a_m^2/(8k^2),&k\ge2.
   \end{cases}
\]
Thus, with $\Pi_0$ the projection onto spatial constants,
\[
 w(T)=\left(\frac{3\pi a_m^2}{8}\Lambda^{-2}
                -\frac{\pi a_m^2}{16}\Pi_0\right)v.
\]
The kernel of $\Lambda^{-2}$ on the unit $S^3$ is
$G(r)=(\pi-r)/(4\pi^2\sin r)$: it solves the homogeneous radial equation
away from its pole, is smooth at the antipode, and has the normalized
singularity $1/(4\pi r)$. Since $\Pi_0$ has kernel $1/(2\pi^2)$, the kernel
of the operator displayed above is
\[
 K(r)=\frac{a_m^2}{32\pi}
       \left(\frac{3(\pi-r)}{\sin r}-1\right)>0\qquad(0<r\le\pi),
\]
using $(\pi-r)/\sin r\ge1$, with its continuous value at $r=\pi$.
Consequently $w(T,Y)>0$. Continuity now permits a choice of a fixed
$s\in(0,\tau)$ with $w(T-s,Y)>0$. Theorem~\ref{thm:diamonds} puts
$x_*=(T-s,Y)$ in $F_0(b)\cap\Isp(b)$, so $\psi_0(x_*)=0$.
Smooth dependence on $\lambda$ yields
\[
 \psi_{m^2}(x_*)=m^2w(x_*)+o(m^2)>0
\]
for every sufficiently small positive $m$. Because $a(x_*)>0$, the physical
field and zero agree on $\Is(b)$ but differ at $x_*$. The event and data
were chosen before the mass interval.
\end{proof}

Theorem~\ref{thm:massive} below gives the full formula; this proposition
supplies an explicit first-order witness. The result concerns a
linear classical field on a fixed background, without interactions,
quantization or gravitational backreaction.
